\section{Sealed Multi-Gate Design}
\label{sec:design}

\subsection{Cohorts and one-shot access}

Table~\ref{tab:gates} summarizes the three confirmatory cohorts.  Each gate has
a development generator, a separate encrypted seed commitment, a zero-access
audit, a frozen source digest, and an atomic one-shot lock.  A sealed result is
not reused to change a threshold, feature, model population, or sample size.

\begin{table}[ht]
\centering
\caption{Confirmatory cohorts.  Runs equal worlds times six controls times
three optimizer seeds.}
\label{tab:gates}
\begin{tabular}{lrrrl}
\toprule
Gate & Worlds & Runs & Max horizon & Primary purpose \\
\midrule
Structural OOD & 50 & 900 & 1 & unseen mechanisms and routing \\
Open-loop rollout & 62 & 1,116 & 32 & recursive structural error \\
Criterion validity & 69 & 1,242 & 8 & plan failure and first failure \\
\bottomrule
\end{tabular}
\end{table}

The OOD and rollout gates use all six controls in paired comparisons.  The
criterion-validity gate fits its hazard model on the four controls that produce
nondegenerate errors.  The dense and oracle-signature controls remain mandatory
structural-zero checks and are never deleted from the report.

\subsection{Structural OOD gate}

Development and sealed mechanisms are disjoint at the world level.  The
generator distinguishes unseen recombinations, unseen structural modes,
unseen mechanism parameters, unseen action compositions, bridge-consistent
shifts, and bridge-violating negative controls.  The primary comparisons pair
correct routing against random and wrong routing while preserving capacity and
compute.  Fifty sealed worlds were selected by the frozen power rule.

\subsection{Open-loop gate}

The rollout protocol separates teacher-forced and recursive predictions at
horizons \(1,2,4,8,16,32\).  It records the area under the
\(\DeltaTheta\) trajectory, the fixed-layer vector, survival without structural
failure, bridge violation, and tracking diagnostics.  The recursive inequality
audit is retained even when its accumulated worst-case numerical bound is
vacuous.

\subsection{Criterion-validity gate}

Each world contains 32 units.  A unit has a domain-separated eight-step probe
and eight tasks with horizons
\[
(1,2,4,8,1,2,4,8).
\]
Task utilities are generated only from structural predicates.  The primary
population contains the layer-routed, strict-factorized, random-routed, and
wrong-routed controls.

The scalar baseline includes training negative log likelihood, one-step and
open-loop latent mean-squared error, endpoint exactness, task horizon and goal
weights, model identity, plan-selection margin, and selected-versus-alternative
scalar error contrasts.  The augmented model adds prespecified task-local
\(\DeltaTheta\), fixed-total, fixed-layer, and goal-aligned structural
contrasts.  A mandatory non-gating sensitivity gives the baseline direct raw
layer mismatches.

A ridge logistic discrete hazard \(h_i(t)\) is fit on development worlds.  Its
cumulative failure probability is
\[
\widehat F_i(t)
=1-\prod_{s=1}^{t}\{1-\widehat h_i(s)\}.
\]
All coefficients, feature order, means, scales, and penalties are frozen before
the sealed seed is revealed.

\subsection{Estimands and inference}

The two co-primary effects are
\[
\begin{split}
\mathcal E_{\mathrm{bin}}
&=
\operatorname{BS}_{\mathrm{scalar}}(8)
-
\operatorname{BS}_{\mathrm{TSI}}(8),\\
\mathcal E_{\mathrm{int}}
&=
\frac{1}{8}\sum_{t=1}^{8}
\left[
\operatorname{BS}_{\mathrm{scalar}}(t)
-
\operatorname{BS}_{\mathrm{TSI}}(t)
\right].
\end{split}
\]
Here, \(\operatorname{BS}(8)\) is the Brier score at terminal time index
\(t=8\) of the fixed eight-step candidate trajectory; the numeral is a time
index, not a task identifier.  The integrated estimand averages the eight time
indices of that trajectory.
Positive values favor the structural augmentation.  Scores are first averaged
within world over units, primary controls, and nested optimizer seeds.
One-sided world-level Student tests use Holm familywise correction
\citep{holm1979simple}.  We also require positive Bonferroni simultaneous
lower bounds from a 20,000-replicate world-cluster bootstrap and from a
world-random-intercept, seed-nested variance decomposition.  Both point effects
must satisfy the prespecified smallest effect size of interest
\(\SESOI=0.0025\).  The inferential null remains a nonpositive effect; we do
not claim that a confidence lower bound exceeds \(\SESOI\).
